% ============================================================
% PASTE-READY SNIPPETS FOR Wikispeedia_Flow_Trees_PLOSOne.tex
% Generated 2026-03-17
% ============================================================
% Each section is labeled with where it goes in the .tex file.
% Search for the placeholder text shown in [[ ]] to find insertion points.
% ============================================================


%% ── ABSTRACT ────────────────────────────────────────────────────────────────
%% Replace everything between \begin{abstract} and \end{abstract}

\begin{abstract}
Understanding how collective navigation behavior emerges from the interaction between
network topology and semantic knowledge is a central challenge in network science and
cognitive science.
We introduce \textit{Flow Trees}---a tree representation that encodes when and how
groups of trajectories diverge through time---and apply them to Wikispeedia, an
online game in which players navigate a curated semantic network of Wikipedia articles
to reach a target.
Using 70,770 human trajectories across 4,592 articles, we construct Flow Trees rooted
at each of the 500 most frequently visited nodes and compare them to matched
random-walker trajectories.
We find that human navigation is more semantically structured than random walks:
node traversal frequency correlates more strongly with lexical generality than with
structural centrality, whereas the reverse holds for random walkers.
Landmark articles---identified by structural out-degree, player traversal frequency,
or lexical generality---produce shallower, more fan-like Flow Trees with higher
branching factors, consistent with a ``zoom-out, home-in'' navigation strategy.
This landmark signature is robust for structural and data-driven definitions but
absent for semantic ones, suggesting that landmark behavior is driven by connectivity
rather than word frequency.
We introduce a Flow-Tree-derived landmark score (branching factor divided by depth)
that identifies dynamically central nodes not captured by structural definitions alone.
Finally, comparing forward and reverse Flow Trees reveals directional asymmetry:
landmark nodes act as stronger convergence points than distributors, indicating that
human navigation produces directed flow structures not present in random walks.
\end{abstract}


%% ── FIGURE 1 CAPTION ────────────────────────────────────────────────────────
%% Replace the current Fig 1 caption ("Simple, high-level schematic.")

\caption{\textbf{Flow Tree construction in Wikispeedia.}
(A) In the Wikispeedia game, players navigate from a randomly assigned start article
to a target article by following hyperlinks within a curated subset of Wikipedia
(4,592 articles, 119,881 directed edges).
An example game (start: \textit{Book}, target: \textit{Ocean}) with a sample
trajectory is shown.
(B) Flow Tree construction.
All trajectories that pass through a given root node are collected.
At each step, trajectories are grouped by their next-visited article; when multiple
distinct choices occur, the group splits and a branching node is created.
The process continues recursively until trajectories terminate or become singletons.
\textit{Forward} Flow Trees encode the structure of outgoing trajectories;
\textit{reverse} Flow Trees are built from time-reversed trajectories and encode
the structure of incoming flow.
}


%% ── FIGURE 2 CAPTION ────────────────────────────────────────────────────────
%% Replace current Fig 2 caption entirely.
%% NOTE: The figure sketch embedded in the draft (labeled "Fig 2: Semantics vs.
%% structure. What drives human behavior?") shows the correlation/centrality
%% analysis -- that is what this caption describes. The LaTeX caption text
%% currently says "Landmark vs. non-landmark Flow Trees", which belongs to
%% Fig 3. Swap accordingly.

\caption{\textbf{Semantic and structural drivers of human navigation.}
(A) Three null models used to isolate human collective knowledge from underlying
graph structure: a uniform random walker (coin-flip), a structural random walker
(next node chosen proportionally to out-degree), and a semantic random walker
(next node chosen proportionally to Zipf word frequency of the article title).
(B) Log-scale scatter plots comparing per-node traversal counts to four structural
centrality measures (degree, betweenness, closeness, PageRank) and one semantic
measure (mean Zipf frequency of article-title words), shown separately for human
players and matched random walkers.
Human traversal counts correlate more strongly with lexical frequency than random
counts do, while random-walker counts correlate more strongly with every structural
centrality measure.
This dissociation indicates that human navigation reflects semantic knowledge
beyond what can be explained by graph structure alone.
}


%% ── FIGURE 3 CAPTION ────────────────────────────────────────────────────────
%% Replace current Fig 3 caption entirely.
%% NOTE: The embedded sketch (labeled "Fig 3: Landmark articles have distinct
%% dynamical signatures") shows the full landmark analysis -- violin plots,
%% FT score scatter, gallery. The current LaTeX caption says "Human vs. random
%% flow", which is wrong. Replace with the text below.

\caption{\textbf{Landmark articles have distinct dynamical signatures.}
(A) Schematic of the conjectured ``zoom-out, home-in'' navigation strategy:
players first navigate to a highly connected hub article, then follow a
semantically guided path toward the target.
(B) Empirical support for the zoom-in hypothesis: distribution of trajectory steps
before and after first arrival at a high-degree hub, showing that players reach
landmark nodes early in their paths.
(C) Overlap among three landmark definitions---data-driven (top 10\% by player
traversal frequency), structural (top 10\% by out-degree), and semantic (top 10\% by
mean Zipf word frequency)---shown as an UpSet plot.
Six articles appear in all three definitions: \textit{United States}, \textit{Water},
\textit{South America}, \textit{London}, \textit{Human}, and \textit{World War II}.
(D) Representative forward Flow Trees for landmark and non-landmark articles under
human and random-walker trajectories.
Landmark Flow Trees are wider and shallower (fan-like), while non-landmark trees
show deeper sequential structure.
(E) Violin plots of five Flow Tree metrics (number of nodes, leaves, depth, diameter,
average branching factor) for landmark vs.\ non-landmark articles, separated by
landmark definition (rows) and trajectory type (human vs.\ random, colors).
Landmark articles defined by structural or data-driven criteria show significantly
higher branching factor and lower depth (see Table~1); semantic landmarks do not
show a consistent pattern across metrics.
(F) Flow-Tree-derived landmark score (mean branching factor / mean depth) plotted
against log out-degree, with points colored by landmark-definition membership.
Articles with high FT scores but low out-degree are dynamically central but
structurally peripheral; a selection of these FT-only nodes is listed.
(G) Gallery of forward Flow Trees for the union of the top-5 articles per definition
(all four definitions), sorted left-to-right by number of definitions claiming the
article, then by FT score.
Border color and membership dots beneath each panel indicate definition membership
(blue: structural; red: data-driven; green: semantic; purple: FT score; amber
border: in two or more definitions).
}


%% ── FIGURE 4 CAPTION ────────────────────────────────────────────────────────
%% Replace current Fig 4 caption.

\caption{\textbf{Forward--reverse asymmetry reveals directional flow structure.}
(A) Three representative node types shown with forward (left) and reverse (right)
Flow Trees.
\textit{Symmetric fans} (top) exhibit high branching factor in both directions,
indicating that trajectories both arrive from and depart to diverse neighbors.
\textit{Symmetric hallways} (middle) show high depth in both directions,
corresponding to articles on extended, stereotyped paths.
\textit{Asymmetric} nodes (bottom) differ markedly between directions, indicating
directional funneling.
(B) Forward (x-axis) vs.\ reverse (y-axis) branching factor and diameter for
each article in the top-500, colored by landmark type (structural, data-driven,
non-landmark).
Points near the diagonal indicate symmetric flow; deviations indicate directionality.
Landmark articles cluster above the diagonal in branching factor: trajectories
converge on landmarks from more diverse directions than they depart from them,
suggesting landmarks act as stronger attractors than distributors.
(C) The same forward--reverse scatter separated by trajectory type (human vs.\
random walker), with 1.5$\sigma$ confidence ellipses.
Random-walker trajectories cluster near the diagonal; human trajectories show
greater asymmetry, particularly for landmark nodes, reflecting the goal-directed
strategy of funneling through hub articles.
}


%% ── IN-TEXT FIXES ───────────────────────────────────────────────────────────

%% 1. RESULTS SECTION OPENER
%% Find: "with high entropy incoming and outgoing. Nice. [add anything else??]"
%% Replace with:

with high entropy incoming and outgoing flow.
We also examine directional asymmetry between forward and reverse Flow Trees,
showing that landmark nodes function as stronger convergence points than
distributors---a signature absent in random walkers.

%% 2. CENTRALITY DEFINITIONS (in "Untangling semantic and structural influences")
%% Find: "degree centrality (what is), betweenness centrality (what is), closeness
%% centrality (what is), and page rank."
%% Replace with:

degree centrality (fraction of total edges incident to each node),
betweenness centrality (fraction of shortest paths passing through each node),
closeness centrality (inverse of the mean shortest-path distance to all other nodes),
and PageRank (stationary probability of a random walk on the graph).

%% 3. BODY TEXT FIGURE REFERENCE FIX
%% Find (line ~154): "See Fig 2 for a visualization of a representative sample."
%% This sentence is in the landmark analysis section -- it refers to representative
%% flow trees. Change to:

See Fig.~3D for a visualization of a representative sample.

%% 4. BODY TEXT FIGURE REFERENCE FIX
%% Find (line ~176): "compare Flow Trees constructed from human trajectories to
%% those generated by random walkers (see Fig 3)"
%% Change to:

compare Flow Trees constructed from human trajectories to those generated
by random walkers (see Fig.~3E)

%% 5. "RABBIT HOLE" SECTION HEADER
%% The section titled "Rabbit hole" (BHV space content) reads as an informal aside.
%% Either remove the heading or rename it:

\subsubsection*{Metric structure of Flow Tree space}

%% 6. ACKNOWLEDGMENTS
%% Find: "We thank Dale Zhou and Aaron[[TODO]] their feedback"
%% Fill in Aaron's last name. Placeholder:

We thank Dale Zhou and Aaron \textbf{[LAST NAME]} for their feedback and
comments on the manuscript.

%% 7. FIG REFERENCE IN DIRECTIONALITY SECTION
%% Find: "See Fig 4A for an example of these three types."
%% Change to (already correct number, just standardize format):

See Fig.~4A for examples of these three types.

%% ── END OF SNIPPETS ─────────────────────────────────────────────────────────
